\documentclass[12pt]{article}
\usepackage[utf8]{inputenc}
\usepackage[T1]{fontenc}
\usepackage{hyperref}
\title{Sio Language Resources for K-12 Curriculum Development}
\author{Compiled by Claude Sonnet 5.5 (claude-sonnet-5-5) for Andr\'as Kornai}
\date{}
\begin{document}
\maketitle

\section{Introduction}
Sio is an Austronesian language of Papua New Guinea. It has about 3,500 speakers on the north coast of the Huon Peninsula in Morobe Province, in the Wasu area. Sources place it in the Western Oceanic group. The first orthography was based on Kate and appeared in a 1953 book of Bible stories, the Small Catechism and hymns, with the missionary Michael Stolz as main translator. The search found little beyond overviews.

\section{Primary Educational Resources}
None were found. No primer, textbook or school programme for Sio appeared in the search results.

\section{Linguistic and Reference Materials}
\begin{itemize}
\item Wikipedia, ``Sio language''. \url{https://en.wikipedia.org/wiki/Sio_language}
\item Omniglot, Sio writing and language page, covering orthography and phonology notes. \url{https://omniglot.com/writing/sio.htm}
\item Glottolog reference record, Reesink 1981, which appeared in the results for Sio. Its exact content was not checked. \url{https://glottolog.org/resource/reference/id/9446}
\item Everyculture, ``Sio Orientation'': an ethnographic overview. \url{https://www.everyculture.com/Oceania/Sio-Orientation.html}
\end{itemize}
No grammar or dictionary was found directly in the search results.

\section{Religious and Community Resources}
The Omniglot page mentions the early mission publications (Bible stories, catechism, hymns), but no online copy was found.

\section{Digital and Audio Resources}
None were found.

\section{Curriculum Development Context}
The language is stable and used by all in its community, but the search found no materials for schools. A curriculum would need new materials built on the existing orthography, with community involvement. Further searches for SIL or Glottolog-listed grammars and word lists would be sensible next steps.
\end{document}